\section{What the Evidence Supports}

\subsection{A completed software path}

The reference implementation now connects raw-data preparation, model fitting,
calibration, explainable decisions, optional guarded updates, accepted-state
recovery, rollback, and bounded streaming operation. Its interface no longer
requires future capture dates for runtime use. The Linux/amd64 image starts
without compiling or downloading and has measured resource headroom in the
agreed local simulation.

The useful software contribution is this connected, inspectable chain and its
failure behaviour. Input-contract validation prevents an incompatible vector
from being treated as comparable traffic. A complete model--threshold pair is
activated atomically. A stopped or restarted process can recover an accepted
detector. Tests show the shared path preserves static and legacy behaviour
while exercising the newer conservative gates.

\subsection{A numerical contribution with a precise boundary}

LM supplies accurate close log-Gamma differences inside the likelihood.
Minka fixed point still estimates alpha, and the complete anomaly score still
includes the multinomial coefficient. Row-wise matrix setup was made cheaper
without substituting a different statistical objective. On the recorded
ordinary-concentration workloads, numerical differences do not alter the
reported decisions; near the multinomial boundary, retained probes demonstrate
a stronger numerical advantage.

Those observations do not establish that LM improves attack recall or is always
faster. Aggregate paper-style timing, production matrix likelihood,
one-window scoring, and constrained full-window processing are different
measurements. Each must keep its data, implementation, and environment attached.

\subsection{Adaptation is a tradeoff, not an automatic improvement}

The legacy threshold method worsens both labelled development comparisons.
The selected legacy refit helps Chromecast but not Samsung relative to static.
Samsung contamination and concentration movement show why convergence and
better in-buffer likelihood are insufficient safety conditions.

The conservative implementation addresses specific failure mechanisms.
Retrospective comparisons support pooled gains under legacy adaptation for the
retained scoring and reference checks. Static and conservative decisions remain
unchanged. At the ten-minute reference settings, Chromecast and Samsung promote
no conservative candidate, so their results do not prove that learning improved
the normal profile. Controlled promotion tests verify
code correctness, while candidate rejections and cross-device tradeoffs remain
part of the empirical findings.

\subsection{Selected scoring and separate temporal research}

The temporal branch establishes why a single nominal date split is not the
same as an audited timestamp timeline. Its rolling folds and episode metrics
show variability hidden by pooled window counts. The retained detector uses
only composition, packet volume, and ARP counts, with fixed references during
updates. Broader category tails fail later specificity checks, while extra
packet-level behaviour lacks a common validated input across the labelled
devices. These choices favour a small count-based method over a collection of
device-specific exceptions.

Temporal folds and incident definitions remain separate evidence. Their
figures must not be mixed with the retained detector comparisons as if one
experiment produced both. Retrospective gains also cannot substitute for
an independently frozen final test.

\subsection{Claims that remain unsupported}

The evidence does not establish a production edge-security product,
physical/ARM suitability, live capture, lossless event delivery, poisoning
immunity, or general performance across arbitrary IoT devices. It also does
not establish an independently locked final result for the current reference
detector, a winning conservative configuration, an optimal window length,
capture-block confidence intervals, or empirical forecasting from the
installed forecast command.

The four-duration study exposes a real tradeoff. Thirty minutes improves
static attack coverage; ten minutes gives earlier alerts and better temporal
localisation. Hourly aggregation delays detection further, while accepted
longer-window threshold updates can substantially increase false alerts.
These results support investigating thirty-minute static operation, not
replacing every mode with one longer duration. They do not establish a universal
optimum or remove calibration-sample and retrospective-evaluation limits.

Resource suitability is measured in the bounded local simulation; physical
embedded targets remain untested. The modest normal-window fits demonstrate
the observed training requirement, without establishing a general
sample-efficiency guarantee. High precision is demonstrated on the
Chromecast and Samsung traces. Guarded adaptation is implemented, but consistent
recovery from legitimate drift has not been demonstrated across devices. Alert explanations
describe statistical evidence, rather than proving that an attack occurred.

These limits are compatible with a defensible thesis. The writing should
answer what was actually measured, explain negative outcomes, and identify the
software and data boundary for each claim. Additional measurements are useful
when they answer a frozen research question; they are not prerequisites for
describing every implemented component.
